\documentclass[11pt,a4paper]{article}

\usepackage[T1]{fontenc}
\usepackage[utf8]{inputenc}
\usepackage{lmodern}
\usepackage{microtype}
\usepackage{geometry}
\geometry{margin=27mm}
\usepackage{amsmath,amssymb,amsthm,mathtools}
\usepackage{enumitem}
\usepackage{booktabs,longtable,array}
\usepackage{xcolor}
\usepackage{hyperref}
\usepackage{url}
\usepackage{fancyhdr}
\usepackage{lastpage}

\hypersetup{
  colorlinks=true,
  linkcolor=blue!55!black,
  citecolor=blue!55!black,
  urlcolor=blue!55!black,
  pdftitle={The Two-Player 3n plus or minus 1 Game - Version 7.0},
  pdfauthor={Grisha Pochuev}
}

\pagestyle{fancy}
\setlength{\headheight}{14pt}
\fancyhf{}
\lhead{Two-Player $3n\pm1$ Game, Version 7.0}
\rhead{Audited status manuscript}
\cfoot{\thepage\ of \pageref{LastPage}}

\newtheorem{theorem}{Theorem}[section]
\newtheorem{proposition}[theorem]{Proposition}
\newtheorem{lemma}[theorem]{Lemma}
\newtheorem{corollary}[theorem]{Corollary}
\theoremstyle{definition}
\newtheorem{definition}[theorem]{Definition}
\newtheorem{conjecture}[theorem]{Conjecture}
\newtheorem{obligation}[theorem]{Open obligation}
\theoremstyle{remark}
\newtheorem{remark}[theorem]{Remark}

\newcommand{\WIN}{\textnormal{\textsc{Win}}}
\newcommand{\LOSS}{\textnormal{\textsc{Loss}}}
\newcommand{\DRAW}{\textnormal{\textsc{Draw}}}
\newcommand{\odd}{\operatorname{odd}}
\newcommand{\Moves}{\operatorname{Moves}}
\newcommand{\asl}{\operatorname{asl}}
\newcommand{\rhoo}{\rho}
\newcommand{\Nat}{\mathbb N}
\newcommand{\Natz}{\mathbb N_0}

\title{\textbf{The Two-Player $3n\pm1$ Game}\\[2mm]
\large An Audited Binary Core, Strengthened Local Normalizations,\\
and the Exact Remaining Global Barrier\\[2mm]
\normalsize Version 7.0}
\author{Grisha Pochuev\\
\texttt{n\_854@mail.ru}\\
\url{https://github.com/Grisha-Pochuev/3n-plus-minus-1-game}}
\date{11 August 2026}

\begin{document}
\maketitle

\noindent
\fbox{\parbox{0.94\textwidth}{
\textbf{Audit status of Version 7.0.}
This document does \emph{not} claim a proof of the all-start no-draw
conjecture.  Earlier versions did make that claim, but hostile review found
that their global marked-router argument used finite states as ranked proof
tokens without first proving the required provenance, and compressed an
attached short-tail module into an unproved exhaustiveness assertion.
Those claims are withdrawn here.  Every unconditional theorem in this
version is proved in the text.  The global conclusion is stated only as a
conditional theorem under four explicit open obligations.  Thus Version 7.0
is intended to be gap-free as a status manuscript, not to conceal an open
step inside a claimed prize solution.
}}

\begin{abstract}
The two-player $3n\pm1$ game is played on positive odd integers.  A move
chooses $3n-1$ or $3n+1$, removes all powers of two, and wins immediately
when the result is $1$.  Arbitrary play can cycle, so the question is whether
optimal play nevertheless has finite remoteness from every start.  We retain
the exact binary conjugation used in earlier manuscripts, in which the two
moves are an expanding map $A$ and a strictly contracting map $B$.  We then
separate four notions that earlier drafts sometimes conflated: a finite
outcome, a routing witness, a retained proof token, and a temporary arithmetic
cursor.

Two local repairs are proved completely.  First, a reverse fixed-tail step
from $Q_r^e(3c)$ either removes one factor of three while preserving an
actual draw, or exposes a genuine loss state and a strictly lower win child
before the draw continues.  Second, the exponent-one predecessor has the
same complete dichotomy; in the exceptional residue classes its two
continuing children form an exact adjacent lift over the actual loss sibling.
We also give a corrected audit of the source-zero valuation-one row and of
the long high-return pair.  These results close several former local gaps,
but they do not establish the global provenance invariant required to
iterate the marked router indefinitely.  The remaining obligations are
isolated precisely, and a complete order-theoretic proof shows that they
would imply the no-draw conjecture.  A companion Python program checks the
arithmetic identities and permanent counterexamples used in the audit; it is
explicitly supporting evidence rather than a proof of the infinite theorem.
\end{abstract}

\tableofcontents

\section{The game, outcomes, and the exact status}

The current state is a positive odd integer $n$.  From $n>1$ the player to
move chooses one of
\[
  3n-1,\qquad 3n+1,
\]
and repeatedly divides the chosen integer by $2$ until it is odd.  Reaching
$1$ wins immediately.  Equivalently, state $1$ is terminal and losing for
the player to move.

For an infinite game graph, the finite-remoteness outcome classes are defined
inductively.  A state is \WIN\ when it has a \LOSS\ child; it is \LOSS\ when
all of its children are \WIN.  These are the least classes closed under the
two rules.  Every remaining state is called \DRAW.

\begin{lemma}[Elementary draw recursion]\label{lem:draw-recursion}
A \DRAW\ state has no \LOSS\ child and has at least one \DRAW\ child.
\end{lemma}

\begin{proof}
A state with a \LOSS\ child is \WIN.  A state all of whose children are
\WIN\ is \LOSS.  Therefore a state in neither finite class has no \LOSS\
child and cannot have all children \WIN; at least one child is \DRAW.
\end{proof}

Arbitrary play does not terminate: in the original game, suitable choices
give $5\to7\to5$.  The target is therefore not termination under all choices.

\begin{conjecture}[No-draw conjecture]\label{conj:main}
Every positive odd starting state has finite remoteness.  Equivalently, the
game has no \DRAW\ positions and optimal play reaches $1$ after finitely
many moves.
\end{conjecture}

Conjecture~\ref{conj:main} is the target of the EUR 500 problem posed by
Ingo Althoefer~\cite{Prize}.  It is \emph{not} promoted to a theorem in this
manuscript.

\section{Binary conjugation}

Write an original odd state as $n=2m+1$ and define
\[
  F(m)=\left\lceil\frac{3m}{2}\right\rceil
       =\frac{3m+(m\bmod2)}{2}.
\]
For later factorization, also put
\[
 J(s)=2F(s)+1=
 \begin{cases}
  3s+1,&s\text{ even},\\
  3s+2,&s\text{ odd}.
 \end{cases}
\]
For a positive integer $x$, let $R(x)$ be the integer obtained by deleting
the maximal alternating suffix of the binary expansion of $x$, and let
$\asl(x)$ denote the number of deleted bits.  Put $R(0)=0$ and
$\asl(0)=0$.
For example,
\[
 R(1001_2)=10_2,
 \qquad
 R(10101_2)=0.
\]

\begin{lemma}[Alternating-suffix factorization]\label{lem:alt-factor}
For every positive integer $z$, put $\delta=1-(z\bmod2)$.  There is an
integer $k\ge1$ such that
\[
  3z+1+\delta=2^k J(R(z)).
\]
If $R(z)>0$, $k$ is the length of the maximal alternating suffix of $z$;
if $R(z)=0$, $k$ is one more than the binary length of $z$.
\end{lemma}

\begin{proof}
First suppose $r=R(z)>0$, let $k=\asl(z)$, and put
$\epsilon=r\bmod2$.  The deleted alternating word has the exact numerical
value
\[
 t_{k,0}=\left\lfloor\frac{2^k}{3}\right\rfloor,
 \qquad
 t_{k,1}=\left\lfloor\frac{2^{k+1}}{3}\right\rfloor,
\]
according as $\epsilon=0$ or $1$, so $z=2^k r+t_{k,\epsilon}$.
Since $J(r)=3r+1+\epsilon$, direct substitution in the two parity cases
gives
\[
  3t_{k,\epsilon}+1+\delta=2^k(1+\epsilon).
\]
Adding the contribution $3\cdot2^k r$ yields the claimed identity.
If $R(z)=0$, the whole binary word is alternating and starts with $1$.
Let $\ell$ be its binary length.  Then $z=t_{\ell,1}$, and the same
calculation gives
\[
  3z+1+\delta=2^{\ell+1}=2^{\ell+1}J(0).
\]
Thus the exponent in the statement is $k=\ell+1$.
\end{proof}

\begin{proposition}[First normal form]\label{prop:first-normal}
For $m>0$, the two children in $m$-coordinates are exactly
\[
  F(m),\qquad F(R(m)).
\]
\end{proposition}

\begin{proof}
Since
\[
  3(2m+1)-1=2(3m+1),\qquad
  3(2m+1)+1=2(3m+2),
\]
exactly one of $3m+1$ and $3m+2$ is odd.  If $m$ is even, the direct odd
expression is $3m+1$ and its new $m$-coordinate is $3m/2=F(m)$.  If $m$
is odd, the direct odd expression is $3m+2$ and its new coordinate is
$(3m+1)/2=F(m)$.

The other expression is
\[
  3m+2-(m\bmod2)=3m+1+\delta,
  \qquad \delta=1-(m\bmod2).
\]
Lemma~\ref{lem:alt-factor} says that its odd part is $J(R(m))$.  Returning
to $m$-coordinates gives
\[
  \frac{J(R(m))-1}{2}=F(R(m)).
\]
Thus the two children are exactly the displayed pair.
\end{proof}

Every child in Proposition~\ref{prop:first-normal} lies in the image of
$F$.  Representing the state $F(q)$ by $q$ gives the conjugated game
\[
  A(q)=F(q)=\left\lceil\frac{3q}{2}\right\rceil,
  \qquad
  B(q)=R(F(q)),
\]
with terminal state $q=0$.

\begin{proposition}[Strict contraction]\label{prop:B-contracts}
For every $q>0$,
\[
  B(q)<q.
\]
\end{proposition}

\begin{proof}
The deletion defining $R$ removes at least one binary digit, so
\[
 B(q)\le \left\lfloor\frac{F(q)}2\right\rfloor
 \le \left\lfloor\frac{3q+1}{4}\right\rfloor<q
\]
for $q\ge2$.  Directly, $B(1)=0$.
\end{proof}

After one original move, the $m$-coordinate is in the image of $F$.
Consequently a proof that every conjugated state $q$ has finite remoteness
would imply Conjecture~\ref{conj:main}; the first original move would then
lead to one of two finite-remoteness states.

\section{Constant tails and coefficient sources}

For a positive odd integer $a$, a phase $e\in\{0,1\}$, and $r\ge1$, define
\[
  Q_r^e(a)=a2^r-e.
\]
This is the unique representation of a positive binary word whose maximal
constant suffix consists of $r$ copies of the bit $e$.

\begin{lemma}[Long-tail recurrence and shared boundary]\label{lem:long-tail}
For every $r\ge3$,
\[
 A(Q_r^e(a))=Q_{r-1}^e(3a),
 \qquad
 B(Q_r^e(a))=Q_{r-2}^e(3a).
\]
Moreover,
\[
 B(Q_1^e(a))=B(Q_2^e(a)).
\]
\end{lemma}

\begin{proof}
The formula for $A$ follows directly from the two phase cases:
\[
 A(a2^r)=3a2^{r-1},
 \qquad
 A(a2^r-1)=3a2^{r-1}-1.
\]
When $r\ge3$, the image still ends in at least two equal bits.  Its maximal
alternating suffix therefore consists only of the last bit, and applying
$R$ removes one further constant bit, proving the formula for $B$.

At the boundary,
\[
 A(Q_1^e(a))=3a-e,
 \qquad
 A(Q_2^e(a))=6a-e=2(3a-e)+e.
\]
The last bit of $3a-e$ is $1-e$, so the appended bit $e$ extends the same
maximal alternating suffix by one digit.  Deleting it leaves the same
prefix in both cases.
\end{proof}

Recall that $J(s)=2A(s)+1$.  The map $J$ is an increasing bijection
from $\Natz$ to the positive odd integers not divisible by $3$: explicitly,
$6u+1=J(2u)$ and $6u+5=J(2u+1)$.  Thus every positive odd coefficient has a
unique factorization
\[
  a=3^k J(s),\qquad k\ge0.
\]
We call $s$ its \emph{coefficient source}.  Every positive state has unique
canonical coordinates
\[
 q=Q_r^e(3^k J(s)).
\]
Indeed, if $q$ is even, take $e=0$, $r=v_2(q)$, and $a=q/2^r$; if $q$ is
odd, take $e=1$, $r=v_2(q+1)$, and $a=(q+1)/2^r$.  In both cases $a$ is
positive and odd, and its unique factorization above determines $k$ and
$s$.  We write $\rho(q)=s$ for this state source.

\begin{lemma}[Ordinary side relation]\label{lem:ordinary-side}
Let $q>0$.  If
\[
 q\not\equiv1,3,12,14\pmod {16},
\]
then
\[
 B(A(q))\in\{A(B(q)),B(B(q))\}.
\]
If $q$ is in one of the four exceptional classes, then $A(q)$ is not in an
exceptional class.
\end{lemma}

\begin{proof}
Put $x=A(q)$, $r=R(x)=B(q)$, and $k=\asl(x)$.  In the
nonexceptional classes one has $r>0$: if $r=0$, then $x$ is wholly
alternating, and the two short words are $10_2,101_2$ while every longer
word ends in $0101$ or $1010$.  Directly, the short words come from
$q=1,3$, and
\[
\begin{aligned}
 A(q)\equiv10\pmod {16}&\Longrightarrow q\equiv1,12\pmod {16},\\
 A(q)\equiv5\pmod {16}&\Longrightarrow q\equiv3,14\pmod {16}.
\end{aligned}
\]
Thus $r=0$ places $q$ in an exceptional class.  A direct residue check gives
\[
\begin{array}{c|c}
 k& q\pmod {16}\\ \hline
 1&0,2,5,7,8,10,13,15\\
 2&4,6,9,11.
\end{array}
\]
Thus every nonexceptional residue has $k\in\{1,2\}$.  We now treat these
two actual suffix lengths, so no fixed modulus is being used to decide an
unbounded suffix.

When $k=1$, the deleted bit equals the last bit
$\epsilon=r\bmod2$, and $x=2r+\epsilon$.  Directly,
\[
 A(x)=2A(r)+\epsilon.
\]
Thus $A(x)$ is obtained from $A(r)$ by appending the bit $\epsilon$.
If this bit equals the last bit of $A(r)$, the maximal alternating suffix
stops at the boundary and $R(A(x))=A(r)$; otherwise it continues through
exactly the maximal alternating suffix of $A(r)$ and
$R(A(x))=R(A(r))=B(r)$.

When $k=2$, the deleted word is $01$ if $r$ is even and $10$ if $r$ is
odd.  Direct substitution gives
\[
 A(x)=
 \begin{cases}
  4A(r)+2,&r\text{ even},\\
  4A(r)+1,&r\text{ odd}.
 \end{cases}
\]
Hence $A(x)$ is obtained from $A(r)$ by appending respectively $10$ or
$01$.  The same boundary test leaves either $A(r)$ or
$R(A(r))=B(r)$.  Since $R(A(x))=B(A(q))$, the first assertion follows.

For completeness, if $q=16u+c$ with
$c\in\{1,3,12,14\}$, then
\[
 A(q)\bmod16\in
 \begin{cases}
  \{2,10\},&c\in\{1,12\},\\
  \{5,13\},&c\in\{3,14\}.
 \end{cases}
\]
None of these residues is exceptional.
\end{proof}

\begin{lemma}[Source-boundary selector]\label{lem:source-selector}
For $s>0$, define
\[
 \alpha(s)=1-\bigl(\lfloor s/2\rfloor\bmod2\bigr).
\]
Factor
\[
 3J(s)+1-2e=2^v J(t),\qquad e\in\{0,1\}.
\]
Then
\[
 (v,t)=
 \begin{cases}
  (1,A(s)),&e=\alpha(s),\\
  (v,B(s))\text{ with }v\ge2,&e=1-\alpha(s).
 \end{cases}
\]
\end{lemma}

\begin{proof}
Put $x=A(s)$.  Its parity is $1-\alpha(s)$.  Since $J(s)=2x+1$,
\[
 3J(s)+1-2e=2(3x+2-e).
\]
For $e=\alpha(s)$ the quantity in parentheses is odd and equals $J(x)$:
it is $3x+1$ when $x$ is even and $3x+2$ when $x$ is odd.  This gives
$v=1$ and $t=A(s)$.

For the opposite phase $e=x\bmod2$, the quantity in parentheses is
\[
 3x+2-(x\bmod2)=3x+1+\bigl(1-(x\bmod2)\bigr).
\]
Lemma~\ref{lem:alt-factor} factors it as a positive power of two times
$J(R(x))=J(B(s))$.  Including the leading factor two gives $v\ge2$ and
$t=B(s)$.
\end{proof}

\begin{lemma}[Exceptional child-pair identity]\label{lem:exceptional-pair}
Let $q\equiv1,3,12,14\pmod {16}$ and put $L=B(q)$.  Then for some
$m\ge1$ and $\delta\in\{0,1\}$, the two children of $A(q)$ are exactly
\[
 \{Q_m^\delta(J(L)),Q_{m+1}^\delta(J(L))\}.
\]
\end{lemma}

\begin{proof}
Write $q=16u+c$.  The four direct formulas are
\[
\begin{array}{c|c|c|c}
 c&L&B(A(q))&A^2(q)\\ \hline
 1&R(6u)&18u+1&36u+3\\
 3&R(6u+1)&18u+4&36u+8\\
 12&R(6u+4)&18u+13&36u+27\\
 14&R(6u+5)&18u+16&36u+32.
\end{array}
\]
For $q=1$, one has $L=0$ and the two children of $A(q)=2$ are
$1=Q_1^1(J(0))$ and $3=Q_2^1(J(0))$, so the assertion is direct.
Otherwise put $z=6u,6u+1,6u+4,6u+5$ in the four rows.  Then $z>0$,
$L=R(z)$, and $B(A(q))=3z+1$.  Lemma~\ref{lem:alt-factor} gives
$3z+1=Q_m^\delta(J(L))$ for some $m\ge1$, where
$\delta=1-(z\bmod2)$.  The last column of the table is
$2B(A(q))+\delta$, hence is $Q_{m+1}^\delta(J(L))$.
\end{proof}

\begin{lemma}[Universal source bound]\label{lem:source-bound}
For every positive state $q$,
\[
  \rho(q)\le\frac{q-1}{6}.
\]
\end{lemma}

\begin{proof}
If $q=Q_r^e(3^k J(s))$, then $r\ge1$ and $e\le1$ give
$q\ge2J(s)-1$.  Since $J(s)\ge3s+1$, we have $q\ge6s+1$.
\end{proof}

\subsection{Two permanent negative regressions}

The first invalid shortcut in earlier manuscripts was
\[
 \text{smaller canonical coefficient}
 \quad\Longrightarrow\quad
 \text{smaller coefficient source}.
\]
It is false in the presence of powers of three.

\begin{proposition}[Coefficient/source separation]\label{prop:coeff-source-counter}
Let $s=1$, so $J(s)=5$, and put $a=3J(s)=15$, $e=1$.  The common boundary
child is
\[
 p=R(3a-e)=R(44)=22.
\]
Its canonical odd coefficient is $11<15$, but $11=J(3)$.  Hence its source
has increased from $1$ to $3$.
\end{proposition}

\begin{proof}
The displayed arithmetic is direct.  In binary, $44=101100_2$ and deleting
its maximal alternating suffix gives $10110_2=22$.  Since
$22=Q_1^0(11)$ and $11=J(3)$, the source is $3$.
\end{proof}

The second invalid shortcut compared a return after several expanding moves
with the numerical source before the streak.

\begin{proposition}[Pre-streak anchor separation]\label{prop:prestreak}
In the conjugated game,
\[
  20\xrightarrow{A}30\xrightarrow{A}45\xrightarrow{A}68,
  \qquad B(68)=25>20.
\]
Thus a contracting return after a multi-$A$ streak need not lie below the
source retained before the streak.
\end{proposition}

\begin{proof}
Each value follows from the definitions of $A$ and $B$.
\end{proof}

Neither implication is used later in this manuscript.

\section{Finite outcomes, witnesses, and retained tokens}

Every finite \WIN\ or \LOSS\ state has a canonical proof height.  Put
$h(0)=0$ and define
\[
 h(x)=1+\min\{h(y):y\text{ is a \LOSS\ child of }x\}
 \quad(x\text{ \WIN}),
\]
\[
 h(x)=1+\max\{h(y):y\text{ is a child of }x\}
 \quad(x\text{ \LOSS}).
\]
A finite state carried only to justify a local parent/child incidence is a
\emph{routing witness}.  A \emph{retained proof token} is a finite state
whose relation to the existing rank has also been proved: it either uses an
explicit one-shot initialization or replaces an already retained token by
proper proof-tree descendants.  Merely discovering that a state is finite
does not make it a retained token.

\begin{lemma}[Multiset token rank]\label{lem:token-rank}
For a finite multiset $M$ of proof heights, let
\[
  \Phi(M)=\mathop{\#}_{h\in M}\omega^h,
\]
where $\#$ is natural ordinal sum.  Replacing one occurrence of height $H$
by finitely many certified descendants of heights strictly below $H$
strictly decreases $\Phi$.
\end{lemma}

\begin{proof}
In Cantor normal form, the coefficient of $\omega^H$ decreases by one.
Every inserted monomial has smaller exponent, so no finite collection of
inserted terms can compensate for that loss.
\end{proof}

\begin{remark}[Why provenance is not cosmetic]
A finite \WIN\ state may have another \WIN\ child whose proof height is
not controlled by the canonical \LOSS\ branch used to prove the parent
\WIN.  Therefore a later finite state cannot be substituted into an
ordinal rank unless it is either initialized by a declared seed or proved to
be a proper descendant of a retained token.  This is the central audit rule
for the remaining sections.
\end{remark}

\section{A complete reverse fixed-tail progress lemma}

The following strengthens the fixed-tail statement used in Versions 5 and
6.  The old statement produced a finite witness but did not prove the
progress needed after that witness appeared.

\begin{lemma}[Reverse fixed-tail progress]\label{lem:fixed-tail-progress}
Let $c$ be positive and odd, $e\in\{0,1\}$, and $r\ge2$.  Put
\[
 X=Q_r^e(3c),\qquad
 G=Q_r^e(c),\qquad
 H=Q_{r+1}^e(c),\qquad
 Y=A(G).
\]
Then
\[
 \Moves(H)=\{X,Y\}.
\]
Let $w=B(X)$.  More precisely,
\[
 w=\begin{cases}
     B(Y),&r=2,\\
     A(Y),&r\ge3.
   \end{cases}
\]
If $X$ is \DRAW, exactly one of the following progress alternatives is
available:
\begin{enumerate}[label=(\roman*)]
\item $H$ is \DRAW; replacing $X$ by $H$ removes one factor of three from
      the displayed coefficient;
\item $H$ is \WIN, $Y$ is \LOSS, $w$ is a \WIN\ child of $Y$ with
      $h(w)\le h(Y)-1$, and the other child $A(X)$ is \DRAW.
\end{enumerate}
In alternative (ii), the continuation is therefore accompanied by the
strict token replacement $Y\mapsto w$ whenever $Y$ has legitimately been
initialized as a token.
\end{lemma}

\begin{proof}
Since $r+1\ge3$, Lemma~\ref{lem:long-tail} gives
\[
 A(H)=Q_r^e(3c)=X.
\]
For $r\ge3$ it also gives
$B(H)=Q_{r-1}^e(3c)=A(G)=Y$; for $r=2$ the shared boundary identity gives
the same equality.  Hence $\Moves(H)=\{X,Y\}$.

For $r=2$, $X=Q_2^e(3c)$ and $Y=Q_1^e(3c)$, so the shared boundary identity
gives $B(X)=B(Y)$.  If $r\ge3$, direct use of the long-tail formulas,
including the elementary $Q_2\mapsto Q_1$ row when $r=3$, gives
\[
 B(X)=Q_{r-2}^e(9c)=A(Y).
\]
This proves the common-child statement.

Assume $X$ is \DRAW.  Since $H$ has a \DRAW\ child, $H$ is not
\LOSS.  If $H$ is \DRAW, alternative (i) holds.  If $H$ is \WIN, one
of its children must be \LOSS.  Since $X$ is \DRAW, that child is
necessarily $Y$, so $Y$ is \LOSS.
The state $w$ is an actual child of this \LOSS\ state and is therefore
\WIN, with $h(w)\le h(Y)-1$.  It is also the $B$-child of $X$.  A \DRAW\
state has no \LOSS\ child and at least one \DRAW\ child.  Since $w$ is
\WIN, the other child $A(X)$ must be \DRAW.  This is alternative (ii).
\end{proof}

\begin{corollary}[Finite reverse factor countdown]\label{cor:reverse-countdown}
Starting from a displayed \DRAW\ state $Q_r^e(3^k a)$ with $r\ge2$ and
$k>0$, repeated use of Lemma~\ref{lem:fixed-tail-progress} either reaches a
factor-free displayed coefficient after finitely many reverse steps, or
exposes a strict finite descendant before continuing the \DRAW.
\end{corollary}

\begin{proof}
Every occurrence of alternative (i) reduces $k$ by one.  The nonnegative
integer $k$ cannot be reduced infinitely.  Alternative (ii) is the stated
strict finite descendant event.
\end{proof}

\begin{remark}
The replacement $X\rightsquigarrow H$ is a reverse proof normalization,
not a legal move from $X$.  Any global use must still preserve an actual
continuing \DRAW\ cursor, as required later in
Definition~\ref{def:complete-router}.
\end{remark}

\section{A complete exponent-one predecessor split}

A reverse factor lemma valid for tail exponent at least two must not be used
at exponent one.  The correct exponent-one predecessor is as follows.

\begin{lemma}[Exponent-one predecessor progress]\label{lem:exp-one-progress}
Let $a$ be positive and odd, $k\ge1$, and put
\[
 X=Q_1^e(3^k a),
 \qquad
 P=Q_2^e(3^{k-1}a).
\]
Then $A(P)=X$.  If $X$ is \DRAW, $P$ is either \DRAW\ or \WIN.

If $P$ is \DRAW, the displayed factor exponent decreases from $k$ to
$k-1$.  If $P$ is \WIN, put $L=B(P)$.  Then $L$ is \LOSS.  Moreover:
\begin{enumerate}[label=(\alph*)]
\item if $P\not\equiv1,3,12,14\pmod{16}$, then
      $w=B(X)=B(A(P))$ is an actual \WIN\ child of $L$,
      $h(w)\le h(L)-1$, and $A(X)$ is \DRAW;
\item if $P\equiv1,3,12,14\pmod{16}$, then the two children of
      $X=A(P)$ form an exact adjacent factor-free pair
      \[
        \{Q_m^\delta(J(L)),Q_{m+1}^\delta(J(L))\}
      \]
      for some $m\ge1$ and $\delta\in\{0,1\}$, and this pair contains a
      continuing \DRAW.
\end{enumerate}
\end{lemma}

\begin{proof}
The formula $A(P)=X$ is the $Q_2\mapsto Q_1$ instance of the constant-tail
calculation.  A state with a \DRAW\ child cannot be \LOSS.  If $P$ is
\WIN, the \DRAW\ child $X$ cannot witness the win, so the other child
$L=B(P)$ is \LOSS.

Assume first that $P$ is ordinary.  Then $L>0$: if $B(P)=0$, the
binary word $A(P)$ is wholly alternating, and the residue argument used in
Lemma~\ref{lem:ordinary-side} puts $P$ in an exceptional class.
Lemma~\ref{lem:ordinary-side} therefore gives
\[
 B(A(P))\in\{A(B(P)),B(B(P))\}=\Moves(L).
\]
Thus $w=B(X)=B(A(P))$ is a child of the \LOSS\ state $L$, hence is \WIN\
and has height at most $h(L)-1$.  Since $X$ is \DRAW\ and its $B$-child
$w$ is \WIN, its other child $A(X)$ is \DRAW.

If $P$ is exceptional, write $P=16t+c$ with
$c\in\{1,3,12,14\}$.  Lemma~\ref{lem:exceptional-pair} gives exactly the adjacent pair over $L=B(P)$.
Because $X$ is \DRAW, neither child is \LOSS\ and at least one is \DRAW.
\end{proof}

\begin{remark}
Lemma~\ref{lem:exp-one-progress} closes the local misuse of reverse factor
removal at exponent one.  It does \emph{not} by itself prove that the newly
exposed $L$ is already a retained token in a global rank.  That is a
separate initialization/provenance question.
\end{remark}

\section{The source-zero valuation-one row, corrected}

Source zero means that the odd coefficient is a power of three.  Put
\[
  S_{k,r}^e=Q_r^e(3^k).
\]
Long tails reduce to $r\in\{1,2\}$ in finitely many forward moves.  The
resulting signed boundary has, for some $n\ge1$,
\[
  3^n+1-2e=2^j J(y).
\]

\begin{lemma}[Exact source-zero valuation table]\label{lem:zero-table}
The valuation is
\[
\begin{array}{c|cc}
 &n\text{ odd}&n\text{ even}\\ \hline
 e=0&j=2&j=1\\
 e=1&j=1&j=2+v_2(n).
\end{array}
\]
When $j=1$,
\[
  y=\frac{3^{n-1}-1}{2}.
\]
\end{lemma}

\begin{proof}
For $e=0$, the lifting-the-exponent formula gives
$v_2(3^n+1)=2$ for odd $n$, while $3^n+1\equiv2\pmod4$ for even $n$.
For $e=1$, $3^n-1$ has valuation $1$ for odd $n$ and
$2+v_2(n)$ for even $n$.  It remains to identify $y$ in the two $j=1$
rows.  Put $y=(3^{n-1}-1)/2$.  If $e=0$, then $n$ is even, so $y$ is odd
and
\[
 J(y)=3y+2=\frac{3^n+1}{2}.
\]
If $e=1$, then $n$ is odd, so $y$ is even and
\[
 J(y)=3y+1=\frac{3^n-1}{2}.
\]
These are exactly the odd quotients after the single factor of two is
removed.
\end{proof}

\begin{lemma}[Correct raw selector in the $j=1$ row]\label{lem:zero-j1}
Assume $j=1$ and put $d=1-e$.  The signed exit is
\[
  T=Q_1^d(J(y)),
\]
and its raw companion is
\[
  R(T)=
  \begin{cases}
    A(y),&e=0,\\
    B(y),&e=1.
  \end{cases}
\]
Suppose this signed/raw pair is exposed as the two children of a \DRAW\
boundary state, the raw companion is \DRAW, and $T$ is finite.  Then $T$ is
\WIN.  The following split is exhaustive:
\begin{enumerate}[label=(\roman*)]
\item if $y$ is \WIN, the other ordinary child of $y$ is an exact
      \LOSS\ sibling of the raw \DRAW;
\item if $y$ is \DRAW, the row yields a reverse \DRAW\ parent $y$, but it
      does not yield a finite token or a source decrease.
\end{enumerate}
\end{lemma}

\begin{proof}
The first identity is the definition of the signed exit.  In the $e=0$
valuation-one row, $n$ is even and
$y=(3^{n-1}-1)/2\equiv1\pmod4$.  Hence $A(y)$ is even and
\[
 T=2J(y)-1=4A(y)+1
\]
is obtained by appending the bits $01$ after a final $0$; deleting that
maximal alternating suffix leaves $A(y)$.  In the $e=1$ row, $n$ is odd and
$y\equiv0\pmod4$, so $A(y)$ is again even and
\[
 T=2J(y)=4A(y)+2.
\]
The appended bits are $10$ and the alternating deletion continues through
exactly the maximal alternating suffix of $A(y)$, leaving
$R(A(y))=B(y)$.

The raw and signed exits are the two children of the same \DRAW\ boundary
configuration.  If the signed exit is finite, it cannot be \LOSS, because a
\DRAW\ state has no \LOSS\ child; hence it is \WIN.  Since the raw exit is
an actual child of $y$, the state $y$ cannot be \LOSS.  If $y$ is \WIN, its
other child must be \LOSS\ because the displayed child is \DRAW.  If $y$
is \DRAW, no finite sibling follows.  These are all possible outcomes of
$y$.
\end{proof}

The second branch of Lemma~\ref{lem:zero-j1} was incorrectly compressed in
Version 6.3 into a generic ``loss-sibling entry.''  It is not such an entry.
It is one of the explicit global obligations below.

\section{High returns: the identities are proved, the rank provenance is not}

The long high-return coordinates used in Version 6.3 are arithmetically
correct.  Let $t>0$, $g\in\{0,1\}$, and $v\ge7$.  Put
\[
 u=Q_{v-1}^g(J(t)),
 \qquad
 c=Q_{v-3}^g(3J(t)),
 \qquad
 p=B(A(u)),
 \qquad
 b=B(p).
\]

\begin{lemma}[Exact long high-return diamond]\label{lem:high-return}
The identities
\[
 p=A(c)=Q_{v-4}^g(9J(t)),
 \qquad
 b=Q_{v-6}^g(27J(t))
\]
hold.  Suppose, in addition, that the outcome guard proves
\[
 u,c,p,b\text{ are \WIN},
 \qquad A(u)\text{ is \LOSS},
 \qquad B(c)\text{ is \LOSS}.
\]
Then
\[
 h(p)\le h(u)-2,
 \qquad
 h(b)\le h(c)-2.
\]
Consequently, if $u,c$ are already retained tokens, the replacement
\[
  (u,c)\longmapsto(p,b)
\]
is strict in the multiset token rank.  A declared one-shot seed may first
install the pair $\{u,c\}$; after that installation the same replacement
is strict, while the initialization itself must be paid by the separate
seed component rather than by $\Phi$.
\end{lemma}

\begin{proof}
Constant-tail reduction gives
\[
 A(u)=Q_{v-2}^g(3J(t)),
 \qquad
 B(A(u))=Q_{v-4}^g(9J(t))=A(c),
\]
where the boundary rows $v=7$ are included by direct substitution.  Applying
the same calculation to $p$ gives the formula for $b$.

Here $B(u)=c$, and the outcome guard says that $c$ is \WIN.  Thus
$A(u)$ is the unique \LOSS\ child of the \WIN\ state $u$.  Since $p$ is a
child of $A(u)$,
\[
 h(p)\le h(A(u))-1=h(u)-2.
\]
Similarly, $v-3\ge4$ makes $c\equiv0$ or $15\pmod {16}$, so $c$ is
ordinary.  Its child $A(c)=p$ is \WIN, while the guard says that $B(c)$ is
\LOSS; hence $B(c)$ is the unique \LOSS\ child witnessing $c$ as \WIN.
Lemma~\ref{lem:ordinary-side} makes $b=B(A(c))$ a child of $B(c)$, and
therefore
\[
 h(b)\le h(B(c))-1=h(c)-2.
\]
Lemma~\ref{lem:token-rank} gives the final conditional conclusion.
\end{proof}

\begin{remark}[The precise old gap]
The hypotheses ``$u$ and $c$ are finite \WIN\ states'' do not imply that
they are already members of a ranked multiset.  Version 6.3 silently made
that promotion.  Lemma~\ref{lem:high-return} deliberately separates the
valid height inequalities from the still missing global initialization and
continuity invariant.
\end{remark}

\section{The two shortest marked tails}

The shortest marked-tail arithmetic can also be stated without overclaim.
Let $C$ be positive and odd and put $b=Q_D^g(C)$.

\begin{lemma}[The $D=1$ row pays an actual loss token]\label{lem:D1}
Assume
\[
 q=B(b)\text{ is \WIN},
 \qquad
 y=A(b)\text{ is \LOSS},
 \qquad D=1.
\]
Factor
\[
 9J(b)+1-2g=2^j J(w).
\]
Then $j\ge2$ and
\[
 w=\begin{cases}
     A(y),&j=2,\\
     B(y),&j\ge3.
   \end{cases}
\]
Thus $w$ is an actual \WIN\ child of the retained \LOSS\ token $y$ and
$h(w)\le h(y)-1$.
\end{lemma}

\begin{proof}
Since $b=2C-g$,
\[
 9J(b)+1-2g=
 \begin{cases}
  2(27C+5),&g=0,\\
  2(27C-5),&g=1.
 \end{cases}
\]
The parenthesized integer is even because $C$ is odd, so $j\ge2$.
With $y=A(b)=3C-g$, the same identity is
\[
 9J(b)+1-2g
 =2\bigl(3J(y)+1-2(1-g)\bigr).
\]
The expression in parentheses is the signed source transition at
$y$ in phase $1-g$.  Lemma~\ref{lem:source-selector} says that valuation one
in that inner transition (total $j=2$) selects $A(y)$, while larger valuation
selects $B(y)$.  Both are children of the
\LOSS\ state $y$.
\end{proof}

\begin{lemma}[Exact $D=2$ lift]\label{lem:D2}
Assume
\[
 q=A(b)\text{ is \WIN},
 \qquad
 y=B(b)\text{ is \LOSS},
 \qquad D=2.
\]
Factor
\[
 9J(b)+1-2g=2^j J(w).
\]
Then $j=1$.  If $y>0$, the returned source is
\[
  w=Q_m^{1-g}(J(y))
\]
for an exact integer $m\ge2$.  If $y=0$, the returned state has coefficient
source zero.
\end{lemma}

\begin{proof}
Now $b=4C-g$, and
\[
 9J(b)+1-2g=
 \begin{cases}
  2(54C+5),&g=0,\\
  2(54C-5),&g=1.
 \end{cases}
\]
The parenthesized integer is odd, so the valuation is exactly one.  Put
$z=3C-g$.  The shared exponent-two/one boundary gives $y=R(z)$.
When $g=0$, the integer $z$ is odd and
Lemma~\ref{lem:alt-factor} gives
\[
  3z+1=2^\lambda J(y).
\]
The returned source is $w=18C+1=2(3z+1)-1$, and therefore
\[
  w=Q_{\lambda+1}^{1}(J(y)).
\]
When $g=1$, the integer $z$ is even and the same lemma gives
\[
  3z+2=2^\lambda J(y).
\]
Now $w=18C-2=2(3z+2)$, so
\[
  w=Q_{\lambda+1}^{0}(J(y)).
\]
Here $\lambda\ge1$ is the maximal alternating-suffix length of $z$ when
$y>0$.  Thus in both phases
$w=Q_{\lambda+1}^{1-g}(J(y))$ and $m=\lambda+1\ge2$.  If $R(z)=0$,
Lemma~\ref{lem:alt-factor} instead gives a constant-tail state over
$J(0)=1$, which is exactly the source-zero alternative.
\end{proof}

Lemma~\ref{lem:D2} proves attachment of the first returned cursor to the
actual \LOSS\ token $y$.  It does not prove that every later raw and signed
exit preserves that attachment until a descendant is paid.  That latter
statement is the second principal global obligation.

\section{The exact remaining global obligations}

The earlier versions attempted to assemble local arithmetic into a marked
transition system with source anchors, one-shot token seeds, proof-token
multisets, and finite control counters.  The order-theoretic idea is sound,
but the semantic refinement to every legal \DRAW\ continuation is not yet
complete.  The missing content is now isolated as follows.

\begin{obligation}[Z: source-zero $j=1$ bootstrap]\label{obl:Z}
In the second branch of Lemma~\ref{lem:zero-j1}, where both the raw child and
its ordinary parent $y$ are \DRAW, construct a marked transition that either
pays a previously retained source/token component or consumes a declared
one-shot bootstrap which cannot be consumed again in the same earlier
fibre.  Prove that every later return to source zero is compatible with this
non-reseeding rule.
\end{obligation}

\begin{obligation}[H: high-return pair provenance]\label{obl:H}
Reserve an inner one-shot seed for the first long high-return pair
$\{u,c\}$, or prove that the pair consists of descendants of already retained
tokens.  After the seed has been consumed, prove for every subsequent long
high return in the same outer fibre that its incoming pair is exactly the
currently retained pair or certified descendants of it.  No merely finite
routing witness may be promoted to the rank.
\end{obligation}

\begin{obligation}[L: complete LOSS-anchored $D=2$ module]\label{obl:L}
Starting from the exact lift of Lemma~\ref{lem:D2}, enumerate every legal
raw/signed outcome branch, every remaining factor exponent, and the terminal
tail classes $m=1,2,3,\ge4$.  Prove that each macro-step either decreases an
explicit finite counter while retaining the same \LOSS\ token $y$, or
replaces $y$ by an actual child or certified descendant before any generic
obligation is re-entered.  In particular, prove that no branch exits as a
fresh free obligation at an unrelated arithmetic cursor.
\end{obligation}

\begin{obligation}[E: semantic exhaustiveness of the combined router]\label{obl:E}
After Obligations~\ref{obl:Z}--\ref{obl:L} are supplied, construct a
marked macro-relation satisfying every clause of
Definition~\ref{def:complete-router} below.  In particular, give a disjoint
and complete case partition for every control transition.  For each
equal-rank cycle, exhibit a natural counter that strictly decreases; for
each control reset, exhibit a preceding strict source or token replacement.
This proof must refine the actual game outcomes, not only an acyclic graph
of declared control labels.
\end{obligation}

\begin{remark}
Lemma~\ref{lem:fixed-tail-progress} closes the former fixed-tail part of the
entry countdown, and Lemma~\ref{lem:exp-one-progress} closes the local
exponent-one misuse.  They reduce Obligation~\ref{obl:E}, but do not imply
Obligations~\ref{obl:Z}--\ref{obl:L}.
\end{remark}

\section{A complete conditional closure theorem}

This section records the exact abstract statement supplied by the marked-rank
architecture.  The theorem is conditional, but the condition is now a
mathematically explicit property rather than the name of a control graph.

Let a marked configuration retain
\[
  \Pi=(s,\eta,\Phi(M),\zeta,\Psi(N),a),
\]
ordered lexicographically.  Here $s,a\in\Natz$ are genuine retained source anchors, $M,N$ are
finite multisets of certified token heights, $\Psi(N)$ is defined by the
same natural-sum formula as $\Phi(M)$, and $\eta,\zeta\in\{1,0\}$ are
one-shot seed bits ordered by $1>0$.  With the usual orders on the natural
and ordinal components, this finite lexicographic product is a well-order.
Temporary sources, phases, valuations, and tail lengths are not components
of $\Pi$ unless an explicit strict comparison installs them as retained
data.

\begin{definition}[Complete marked-router hypothesis]
\label{def:complete-router}
A marked macro-relation $\Longrightarrow$ is called complete for this game
when all of the following hold.
\begin{enumerate}[label=(C\arabic*)]
\item Every conjugated \DRAW\ state produces an initial marked configuration
      carrying that same actual \DRAW\ as its continuing cursor.
\item From every marked configuration carrying an actual \DRAW, a finite
      outcome-compatible normalization produces a successor marked
      configuration carrying an actual \DRAW.  Reverse-parent transformations
      are permitted only as proof transformations; whenever the cursor follows
      play, it follows a legal \DRAW\ child.
\item Every macro-step weakly decreases $\Pi$.  A token-changing step is a
      one-shot seed or a certified proper-descendant replacement; a
      source-changing step uses an explicit strict comparison of retained
      sources.
\item For each fixed value of $\Pi$, the restriction of
      $\Longrightarrow$ to that fibre is well-founded.  In particular, every
      equal-$\Pi$ control cycle has an explicitly decreasing natural counter.
\end{enumerate}
\end{definition}

\begin{lemma}[Well-founded fibre principle]\label{lem:fibre}
Let a transition relation carry a value in a well-order.  Suppose every
transition weakly decreases that value and the restriction of the relation
to each fixed-value fibre is well-founded.  Then the full relation is
well-founded.
\end{lemma}

\begin{proof}
An infinite path would induce a nonincreasing sequence in a well-order.  If
it had infinitely many strict decreases, those decreases would form an
infinite descending subsequence, which is impossible.  Hence its tail would
lie in one fixed fibre, contradicting fibre well-foundedness.
\end{proof}

\begin{theorem}[Conditional no-draw closure]\label{thm:conditional}
If a complete marked macro-relation in the sense of
Definition~\ref{def:complete-router} is constructed from the game relation,
then the conjugated game has no \DRAW\ positions.  Consequently every
original positive odd start has finite remoteness.
\end{theorem}

\begin{proof}
Assume a conjugated \DRAW\ exists.  Clause (C1) gives an initial marked
configuration carrying it.  Clause (C2), applied recursively, gives an
infinite sequence of marked configurations, each carrying an actual
continuing \DRAW.  Clauses (C3)--(C4) and
Lemma~\ref{lem:fibre} say that the same macro-relation is well-founded, so
such an infinite sequence cannot exist.  Hence no conjugated \DRAW\ exists.
Proposition~\ref{prop:first-normal} transfers the conclusion to every
original odd start: after the first move both possible children are
conjugated finite-remoteness states, so the original state is itself
\WIN\ or \LOSS\ with finite height.
\end{proof}

\begin{remark}
Theorem~\ref{thm:conditional} is not a proof of
Conjecture~\ref{conj:main}.  Obligations~\ref{obl:Z}--\ref{obl:L} are
specific missing semantic modules, and Obligation~\ref{obl:E} is the
remaining task of constructing a relation satisfying
Definition~\ref{def:complete-router}.  None is assumed silently in an
unconditional theorem.
\end{remark}

\section{Computational verification boundary}

The companion program \texttt{verify\_v7\_0.py} checks, over configurable
finite ranges:
\begin{enumerate}[label=(\arabic*)]
\item the original-to-binary normal form and strict contraction of $B$;
\item the long-tail recurrence and shared boundary;
\item the alternating-suffix factorization, the bijection $J$, ordinary
      coefficient-source recovery, the ordinary side relation, the
      source-boundary selector, and the exceptional child-pair identity;
\item the two permanent negative regressions;
\item every arithmetic identity in Lemmas
      \ref{lem:fixed-tail-progress}, \ref{lem:exp-one-progress},
      \ref{lem:zero-table}, \ref{lem:zero-j1},
      \ref{lem:high-return}, \ref{lem:D1}, and \ref{lem:D2}.
\end{enumerate}
The program does not infer infinite outcome classes and does not check
Obligations~\ref{obl:Z}--\ref{obl:E}.  Its final message explicitly says
that the run is supporting evidence, not a no-draw proof.

\section{Hostile-audit matrix}

\renewcommand{\arraystretch}{1.18}
\begin{longtable}{>{\raggedright\arraybackslash}p{0.24\textwidth}
                  >{\raggedright\arraybackslash}p{0.17\textwidth}
                  >{\raggedright\arraybackslash}p{0.49\textwidth}}
\toprule
Item & Status in v7.0 & Audit conclusion \\
\midrule
\endfirsthead
\toprule
Item & Status in v7.0 & Audit conclusion \\
\midrule
\endhead
Binary conjugation and $B(q)<q$ & Proved & Exact algebra; finite regression included. \\
Long-tail identities & Proved & Uniform in $a,r,e$; no bounded-modulus assumption. \\
Coefficient versus source & Corrected & The explicit $1\to3$ counterexample is retained permanently. \\
Multi-$A$ numerical anchor & Corrected & The $20\to30\to45\to68$, $B(68)=25$ counterexample is retained permanently. \\
Fixed-tail reverse entry & Proved locally & Lemma~\ref{lem:fixed-tail-progress} adds the missing continuing-DRAW and descendant-token alternatives. \\
Exponent-one reverse entry & Proved locally & Lemma~\ref{lem:exp-one-progress}; no use of the $r\ge2$ reverse lemma at $r=1$. \\
Source-zero $j=1$ row & Partly proved & Arithmetic and outcome split proved; global non-reseeding bootstrap is Obligation~\ref{obl:Z}. \\
Long high-return heights & Proved conditionally on token legitimacy & The inequalities are valid; rank insertion is isolated as Obligation~\ref{obl:H}. \\
Marked tail $D=1$ & Proved & Actual child of the retained LOSS token. \\
Marked tail $D=2$ & Exact first lift proved & Full raw/signed provenance remains Obligation~\ref{obl:L}. \\
Complete router & Open for the target theorem & No unconditional exhaustiveness claim; Obligation~\ref{obl:E}. \\
No-draw conclusion & Conditional only & Theorem~\ref{thm:conditional} under
Definition~\ref{def:complete-router}; Conjecture~\ref{conj:main} remains
open here. \\
\bottomrule
\end{longtable}

\section{Conclusion}

Version 7.0 does not repeat the cycle in which one local repair is followed
by a new hidden global assumption.  It makes three changes of principle.
First, it proves complete local progress statements for the fixed-tail and
exponent-one reverse rows rather than merely naming a finite witness.
Second, it never identifies a routing witness with a retained token without
a seed or a descendant comparison.  Third, it records the prize statement
as a conjecture until the source-zero bootstrap, high-return pair continuity,
LOSS-anchored $D=2$ module, and combined semantic exhaustiveness have all
been proved.

The result is a manuscript whose stated theorems are auditable and whose
remaining barrier is explicit.  It is not yet a solution of the prize
problem.  Any future Version 8 claiming the no-draw theorem must contain
proofs of Obligations~\ref{obl:Z}--\ref{obl:E}, or replace the marked-router
architecture by a genuinely different complete argument.

\section*{Acknowledgment of the audit}
Questions from Ingo Althoefer and subsequent hostile review exposed the
coefficient/source error, the exponent-one misuse, and the later token
provenance failures.  Version 7.0 treats those findings as permanent
regressions rather than cosmetic edits.

\begingroup
\raggedright
\begin{thebibliography}{9}
\bibitem{Guy1986}
R. K. Guy,
``John Isbell's Game of Beanstalk and John Conway's Game of Beans-Don't-Talk,''
\emph{Mathematics Magazine} 59(5) (1986), 259--269.
\href{https://doi.org/10.1080/0025570X.1986.11977258}{doi:10.1080/0025570X.1986.11977258}.

\bibitem{Guy1996}
R. K. Guy,
``Unsolved Problems in Combinatorial Games,'' Problem 42,
in \emph{Games of No Chance}, MSRI Publications 29,
Cambridge University Press, 1996.

\bibitem{AHZ2024}
I. Althoefer, M. Hartisch, and T. Zipproth,
``Analysis of a Collatz Game and Other Variants of the $3n+1$ Problem,''
in \emph{Advances in Computer Games}, LNCS 14528 (2024), 123--132.
\href{https://doi.org/10.1007/978-3-031-54968-7_11}{doi:10.1007/978-3-031-54968-7\_11}.

\bibitem{Prize}
I. Althoefer,
``Collatz Problems: The Two-Player $3n{+}{-}1$ Game,''
\url{https://althofer.de/collatz-prizes.html}, accessed 11 August 2026.

\bibitem{Zenodo}
G. Pochuev,
``Optimal Play in the Two-Player $3n\pm1$ Game,''
Zenodo, 2026.
\href{https://doi.org/10.5281/zenodo.21844684}{doi:10.5281/zenodo.21844684}.
The archived record predates the hostile audits summarized here.

\bibitem{Repo}
G. Pochuev,
``Optimal $3n\pm1$ Game: Proof and Verification Repository,'' 2026.
\url{https://github.com/Grisha-Pochuev/3n-plus-minus-1-game}.
\end{thebibliography}
\endgroup

\end{document}
